\documentclass[11pt]{article}
\usepackage[margin=1in]{geometry}
\usepackage{amsmath,amssymb,amsthm,mathtools}
\usepackage{microtype}
\usepackage{xcolor}
\usepackage[hidelinks]{hyperref}
\usepackage{enumitem}
\setlist{nosep}
\newtheorem{theorem}{Theorem}
\newtheorem{proposition}{Proposition}
\newtheorem{lemma}{Lemma}
\newcommand{\R}{\mathbb{R}}
\newcommand{\id}{\mathrm{id}}
\title{An Exact Obstruction to the Local-Finiteness Premise\\in Fu's Theorem 2.8}
\author{}
\date{}

\begin{document}
\maketitle

\begin{abstract}
We give an explicit allowed polynomial system for which the regular-parameter
set appearing in Theorem 2.8 of Y.~Fu, \emph{Towards Trans-Exponential
O-minimal Expansion of the Real Ordered Field}, cannot possess the locally
finite closed-simplex triangulation required by the argument.  The example is
the one-variable system $P(x,l_1,l_2,\epsilon,\delta)=x^2$.  On a fixed
identity-coordinate slice, the regular-parameter set is exactly
$[-1,1]\setminus\{0\}$, hence is not closed.  A locally finite union of closed
simplices in the full ambient parameter space is closed, giving an immediate
contradiction.  If the word ``simplex'' is instead interpreted as an open
simplex interior, the membership conclusion used in Claim 2.9 does not follow
from density and local finiteness.  Thus the displayed construction route has
a precise gap.  This note does not claim a counterexample to o-minimality or a
solution of the general trans-exponential problem.
\end{abstract}

\section{The source premise}

For a system $P$ of the form used in Fu's equation (3), equations (4) and (6)
define a parameter set $G$ and the regular locus
\[
 R=\bigl\{(A,\eta,l,\epsilon,\delta)\in
 G\times[-1,1]^{2n+2}: P\circ A-\eta
 \text{ has nonsingular zeros only}\bigr\}.
\]
Theorem 2.8 assumes a triangulation $\{\Delta_\alpha\}$ of $R$ with the
following ambient local-finiteness property: every point of
$G\times[-1,1]^{2n+2}$ has a neighborhood intersecting only finitely many
$\Delta_\alpha$.  The proof then uses density of $R$ to assert in Claim 2.9
that each point of the compact identity slice belongs to one of finitely many
such simplices.

\section{The explicit system}

Take
\[
 n=1,\qquad k=u=0,\qquad
 P(x,l_1,l_2,\epsilon,\delta)=x^2.
\]
Then $m=3n+k+u+2=5$.  Fix arbitrary
$(l,\epsilon,\delta)\in[-1,1]^4$ and define
\[
 p_\eta=(I_5,\eta,l,\epsilon,\delta),\qquad -1\leq\eta\leq1.
\]

\begin{lemma}
For every $\eta\in[-1,1]$, the pair $(I_5,\eta)$ lies in Fu's set $G$.
\end{lemma}

\begin{proof}
The zero set of $x^2-\eta$ is empty for $\eta<0$, consists of the single
$x$-coordinate $0$ for $\eta=0$, and consists of the two $x$-coordinates
$\pm\sqrt\eta$ for $\eta>0$.  Including the four auxiliary coordinates gives
respectively the empty set, $\{0\}\times\R^4$, or
$\{-\sqrt\eta,\sqrt\eta\}\times\R^4$.  For any auxiliary point, restrict to a
bounded neighborhood in those four coordinates.  The inverse image under the
auxiliary-coordinate projection is then empty, one bounded copy, or two
bounded copies.  This is exactly the local boundedness condition defining
$G$.
\end{proof}

\begin{proposition}
On the displayed identity-coordinate slice,
\[
 \{\eta\in[-1,1]:p_\eta\in R\}=[-1,1]\setminus\{0\}.
\]
\end{proposition}

\begin{proof}
For $\eta<0$, the equation $x^2-\eta=0$ has no real zeros, so the assertion
that all zeros are nonsingular is vacuously true.  For $\eta>0$, the only
zeros are $\pm\sqrt\eta$ and the derivative $2x$ is nonzero at both.  At
$\eta=0$, the zero $x=0$ has derivative zero and is singular.  These three
cases prove the equality.
\end{proof}

\section{Failure of ambient local finiteness}

\begin{theorem}
The regular locus $R$ for the system above admits no triangulation by closed
simplices that is locally finite in the full ambient space
$G\times[-1,1]^4$.
\end{theorem}

\begin{proof}
Assume that $R=\bigcup_\alpha\Delta_\alpha$, where each $\Delta_\alpha$ is a
closed simplex and the family is locally finite in the full ambient space.
For $j\geq1$, put $\eta_j=-1/j$.  Then $p_{\eta_j}\in R$ and
$p_{\eta_j}\to p_0$, while $p_0\notin R$.

Choose an ambient neighborhood $U$ of $p_0$ that meets only finitely many
simplices, say $\Delta_1,\ldots,\Delta_N$.  All sufficiently large
$p_{\eta_j}$ lie in $U$, so infinitely many of them belong to one fixed
$\Delta_i$.  Because $\Delta_i$ is closed, their limit $p_0$ also lies in
$\Delta_i\subset R$, a contradiction.
\end{proof}

Equivalently, a locally finite union of closed subsets is closed.  The
sequence above exhibits directly that $R$ is not closed in the ambient space,
so the first bullet of Theorem 2.8 cannot hold under the closed-simplex
meaning used by its proof.

\section{Why open simplex interiors do not repair the argument}

One might instead call the relative interiors of simplices the members of a
triangulation.  Such pieces need not be closed, so the preceding topological
contradiction would not apply verbatim.  However, Fu's Claim 2.9 uses density
and local finiteness to conclude that a limit point on the identity slice
\emph{belongs} to one of the finitely many simplices.  For a decomposition by
open simplex interiors, a sequence may converge to a proper face without the
limit belonging to any of the interiors containing the sequence.  The
displayed points $p_{-1/j}\to p_0$ give exactly this failure mode.  Thus either
the simplices are closed, in which case the assumed ambient local finiteness
is impossible, or they are open interiors, in which case Claim 2.9 requires a
new argument.

\section{Formal and independent verification}

The accompanying Lean 4 development proves:
\begin{enumerate}
  \item the unrestricted regular-parameter set for $x^2-\eta$ is
  $\R\setminus\{0\}$;
  \item the constrained slice is $[-1,1]\setminus\{0\}$;
  \item this slice is not closed; and
  \item it cannot be a locally finite union of closed pieces.
\end{enumerate}
The development compiles with Lean 4.31.0, contains no proof holes, and passes
the recorded axiom audit.  A separate Reviewer checked the source definitions,
the admissibility of the $P=x^2$ system, the mapping from the paper's parameter
space to the formal slice, and the use of the closed-simplex interpretation in
Claim 2.9.

\section{Scope of the conclusion}

The conclusion is a correction to the proof route, not a counterexample to
the conditional conclusion of Theorem 2.8 and not a counterexample to
o-minimality.  The example proves that the stated triangulation premise is not
available for all allowed systems and that the proof cannot obtain it merely
from density of the regular locus.  Consequently, the o-minimality construction
depending on this step requires a replacement argument or a strengthened
hypothesis.

\section*{Primary source}

Yayi Fu, \emph{Towards Trans-Exponential O-minimal Expansion of the Real
Ordered Field}, arXiv:2604.03477v2.  Relevant locations: equations (3), (4),
and (6); Lemmas 2.5--2.7; Theorem 2.8; and Claim 2.9.

\end{document}
